%*****************************************
\chapter{Future Work}\label{ch:futurework}
%*****************************************
The evaluation and analysis of the methodology follow the structure of the approach, separated into \ac{lcw} and \ac{rag} and a \enquote{Common} Section.
Both method discussions are further subdivided into the individual pipeline steps.
Here, however, at the macro level.
%
\section{Common}

The selection of the \ac{lm} is one of the decisive factors for the performance of a \ac{qa}-system such as the one presented here.
The search for the best application-specific models can therefore be significantly expanded.
For \ac{qa} applications with, for example, extensive scientific text, current \acp{llm} can be selected.
These \ac{sota} models have context windows that no longer limit the input of one or more knowledge bases due to their content size.
In addition, they enable multimodal inputs.
This would allow \cite{bb2} to be parsed into the model as context in its unedited form, for example, in PDF format.
However, as already described, the implementation of these \ac{lm} in individual applications is not without limitations.
First, such a model must be available for public access outside generative \ac{ai} platforms, such as \textit{ChatGPT} or \textit{Gemini}.
Access can be gained either via an aggregator or by running the model on local hardware.

%
\subsection{Prompt Engineering}\label{sub:prompt_engineering}
As mentioned in the approach, prompt engineering would significantly improve the generated output.
By using current prompting techniques, it would be possible to guide the respective \ac{lm} to the ideal answer.
As implemented in the \textit{Blablador} pipeline, a clear restriction on the output length within the prompt avoids truncating generated information by a hard-coded \texttt{max\_token} parameters.
Such a small adaption has already a significant impact on the output quality.

The \textit{llm-as-a-judge} is an efficient feature, but its quality and reliability of output also depend on the instructing prompt.
Such an approach could improve the output with regard to efficiency and quality.

The choice of the \ac{lm} provider influences the usability, availability and capability of the respective \ac{lm}.
 %
\section[Improvements - LCW pipeline]{Potential Improvements for the LCW pipeline}

Future work could focus on various structural and operational improvements to the \ac{lcw} pipeline.
As this pipeline is straightforward without much complexity, its main architecture of monolithic context-capture presents challenges in terms of cost, latency and noise sensitivity.
\subsection{Query Normalisation}
As identified in the evaluation, the performance of the \ac{lcw} pipeline is negatively affected when queries contain spelling errors.
The direct comparison revealed significant weaknesses in terms of the models’ robustness.
For a scientific study such as this one, it is of interest to process the deliberately \enquote{incorrect} input data without modification.
As it represents the realistic human input.

To address this weakness, future work should incorporate an automated step to normalise the queries.
This is already standard practice in commercial \ac{ai} chat platforms, such as \textit{ChatGPT} and \textit{Gemini}.
In these systems pre-processing (query rewriting, spelling correction, optimised prompting) of the initial input runs automatically in the background.
Only in this way will a realistic assessment of the model’s performance, as well as the full utilisation of its actual capabilities, be meaningful.
This approach complements \cref{sub:prompt_engineering}.

\subsection{Context Structuring}
Currently, the knowledge base is read and parsed directly into the prompt template in an unprocessed plain text format.
However, there are no clear structural boundaries.
This can lead to instructions and context becoming blurred, which confuses the model and potentially causes it to lose its train of thought.
Future implementations could be adapted to incorporate structured context integration.
Semantic XML markers (e.g. by extending the explicit separation of chapters, sections or definitions) and embedding structural guides could be used to divide the text.
Thus, the model’s internal attention mechanisms could locate and refer to relevant facts more precisely.
\section[Improvements - RAG Pipeline]{Potential Improvements for the RAG Pipeline}

Future work may investigate the optimisation of the \ac{rag} pipeline.
The implementation within this thesis was designed to be functional and effective.
It has the primary goal of a comparative evaluation, rather than being fine-tuned for peak performance. 
Consequently, several parameters, such as \enquote{top\_k} and \enquote{Chunk\_size}, were set to sensible, commonly used values without empirical \textit{A/B testing}.
Systematically and detailed testing of the configuration settings could significantly enhance the efficiency and the quality of its output.
%
\subsection{Chunk Size}
Possible improvements start with the configuration of the embedding.
The size of the chunks should be adapted to the source.
An in-depth analysis of the complexity and style of the literature could help to define the ideal chunk size and the appropriate overlap.
Detailed literature with a high number of concrete, clearly formulated facts per sentence may allow for a smaller chunk size and lower overlap.
On the other hand, a text with extensive wording and paraphrasing of complex content requires larger chunks and overlap to not lose the context.
%
\subsection{Top-K}
Further optimisation of \ac{rag} can be achieved by adjusting the parameter \enquote{top-k}, the number of document snippets retrieved during the retrieval process.
The choice of \textit{k} is a critical consideration.
Retrieving a larger number of chunks extracted from the database may seem advantageous, as they provide more context to the \ac{llm}.
However, this also increases the likelihood that irrelevant or redundant information (\enquote{noise}) will be processed.
This \enquote{noise} can dilute the relevant context.
Thereby, possibly causing relevant details to be lost from focus and leading to a deterioration in the \textit{precision}.
An extending experimental study would need to vary the parameter value while analysing the quality of the generated output by the respective \ac{llm}.
Identifying the highest possible quality of output likely correlates with the inclusion of a \textit{k} and a \textit{threshold}.
%
\subsection{Cost}
Furthermore, practical considerations of cost and latency are directly tied to the number and size of the retrieved chunks. 
A higher \textit{k} value results in a longer prompt being sent to the generative model.
This then leads to higher \ac{api} costs and longer response times from the \ac{lm} provider.
For the target group of students and practitioners operating under resource constraints, these factors are of significant importance.

Future research should therefore aim to identify the optimal balance.
A top-k value between $5$ and $10$ chunks provides a good compromise for many applications. 
A higher value ($10\geq$) should probably only be considered when it is known that the required information is highly distributed throughout the knowledge base.
%

In order to go beyond heuristics, a systematic investigation is recommended. 
This could involve conducting a series of A/B tests for different parameter settings, different implemented components within the respective method pipelines.

Furthermore, parameter settings and adaptation, as well as retrieval and filter techniques exist to improve the generation output.
The following improvements and adaptations could be implemented:
\begin{itemize}
    \item A more capable embedding model, leading to higher vector dimensions and possibly improving context consistency and answer generation. 
    \item Semantic attribute weighting for the retrieval.
    Prioritising certain metadata and corresponding chunks depending on the knowledge base and the current topic of question.
    \item Maximal Marginal Relevance (MMR). Apply MMR for document re-ranking, to ensure the documents sent to the generator are both relevant and diverse.
\end{itemize}
The possible improvements could be investigated using an empirical approach.
This would allow for the ideal configuration for a specific knowledge base and question set to be determined.
The result would be a valuable blueprint for developing a highly optimised, rather than functional, \ac{qa} system.

\section{Security}
Aggregators for \ac{sota} models such as the Jülich Research center (FZJ) and its \textit{Helmholtz AI} scientific platform emphasise the significant role of \textbf{Data Privacy and Security} in \ac{ai} applications.
In \ac{ai} application users often share data without restriction and without critical consideration.
This aspect has become even more important in the era of \textit{agentic systems}.
\ac{lm} in the role of an agent are allowed to act unrestricted and autonomously on private systems and data.
The data is then processed in the black box of the underlying platform processes, uncontrolled and unregulated by the user.
\textit{Helmholtz AI} guarantees a transparent data privacy, but that is not a common standard on commercial \ac{ai} platforms and \ac{lm} aggregators.